\documentclass[12pt, a4paper]{article}

% ============================================================
% PAQUETES
% ============================================================
\usepackage[utf8]{inputenc}
\usepackage[T1]{fontenc}
\usepackage[spanish]{babel}
\usepackage{amsmath, amssymb, amsthm}
\usepackage{mathtools}
\usepackage{geometry}
\usepackage{booktabs}
\usepackage{hyperref}
\usepackage{xcolor}
\usepackage{lmodern}
\usepackage{parskip}
\usepackage{titlesec}
\usepackage{enumitem}
\usepackage{mdframed}

\geometry{top=2.5cm, bottom=2.5cm, left=3cm, right=3cm}

% Colores institucionales
\definecolor{bayesblue}{RGB}{0, 70, 127}
\definecolor{ecogreen}{RGB}{34, 100, 56}
\definecolor{humanred}{RGB}{150, 30, 30}
\definecolor{ghostgray}{RGB}{80, 80, 80}
\definecolor{blockbg}{RGB}{245, 248, 252}

% Formato de títulos
\titleformat{\section}{\large\bfseries\color{bayesblue}}{}{0em}{\thesection.\quad}
\titleformat{\subsection}{\normalsize\bfseries\color{bayesblue!80}}{}{0em}{\thesubsection.\quad}

% Entorno para notas técnicas
\newmdenv[
  backgroundcolor=blockbg,
  linecolor=bayesblue!40,
  linewidth=1pt,
  leftmargin=0pt,
  rightmargin=0pt,
  innerleftmargin=10pt,
  innerrightmargin=10pt,
  innertopmargin=8pt,
  innerbottommargin=8pt,
  skipabove=8pt,
  skipbelow=8pt
]{notabox}

% ============================================================
\begin{document}

% ============================================================
% PORTADA
% ============================================================
\begin{center}
  {\LARGE\bfseries\color{bayesblue} Lógica del Modelo Bayesiano de Ocupación}\\[0.5em]
  {\large Construcción de la Posterior para el Copetón (\textit{Zonotrichia capensis})}\\[0.3em]
  {\small Proyecto Estadística Bayesiana --- Bogotá, 2025}
\end{center}

\vspace{0.5em}
\noindent Esta guía detalla el proceso matemático y lógico para construir un modelo Bayesiano de ocupación jerárquico, enfocado en el cruce de datos de avistamientos (eBird) y contaminantes atmosféricos ($\text{PM}_{10}$, $\text{O}_3$) en Bogotá.

\hrule
\vspace{1em}

% ============================================================
\section{El Fundamento: El Teorema de Bayes}
% ============================================================

La meta de cualquier análisis Bayesiano es hallar la \textbf{Distribución Posterior} de los parámetros ($\theta$), dado que hemos observado ciertos datos ($D$):

\begin{equation}
  P(\theta \mid D) = \frac{P(D \mid \theta)\, P(\theta)}{P(D)}
\end{equation}

Donde cada término cumple un rol puntual:

\begin{itemize}
  \item $P(D \mid \theta)$ es la \textbf{Likelihood} (Verosimilitud): ¿Qué tan probables son los datos según mis parámetros?
  \item $P(\theta)$ es la \textbf{Prior}: Mi conocimiento o suposición inicial sobre los parámetros.
  \item $P(D)$ es la \textit{evidencia}, una constante de normalización que suele ignorarse en el muestreo MCMC.
\end{itemize}

\hrule
\vspace{1em}

% ============================================================
\section{Especificación de las Priors (No Informativas)}
% ============================================================

Para cada coeficiente de las regresiones, se asigna una distribución previa asumiendo ignorancia completa, pero permitiendo que el muestreo converja y que los datos sean los únicos árbitros de la respuesta.

\subsection{Estructura y Variables Específicas}

Según el análisis exploratorio de datos (EDA), los hiperparámetros se configuran sobre dos componentes aislados:

\begin{itemize}
  \item \textbf{\textcolor{ecogreen}{Ecuación de Ocupación Ecológica ($\psi$):}} Estructurada por $\beta_0$ (intercepto base), $\beta_{\text{PM10}}$ (efecto del $\text{PM}_{10}$) y $\beta_{\text{O3}}$ (efecto del Ozono).
  \item \textbf{\textcolor{humanred}{Ecuación de Detección/Esfuerzo ($p$):}} Estructurada por $\alpha_0$ (intercepto de detectabilidad), $\alpha_{\text{duracion}}$ (tiempo en minutos), $\alpha_{\text{estacion}}$ (heterogeneidad geográfica) y $\alpha_{\text{protocolo}}$.
\end{itemize}

\subsection{Distribuciones a Priori}

Para acoplarse al aumento de variables de Pólya-Gamma y a la función logística (Logit), todas las pre-creencias sobre los coeficientes se modelan con \textbf{Distribuciones Normales} con hiperparámetros estandarizados no informativos. Cada coeficiente distribuye de manera independiente:

\medskip
\textbf{\textcolor{ecogreen}{Eje Ecológico ($\psi$):}}
\begin{align}
  \beta_0 &\sim \mathcal{N}(0,\, 1) \\
  \beta_{\text{PM10}} &\sim \mathcal{N}(0,\, 1) \\
  \beta_{\text{O3}} &\sim \mathcal{N}(0,\, 1)
\end{align}

\medskip
\textbf{\textcolor{humanred}{Eje de Detección ($p$):}}
\begin{align}
  \alpha_0 &\sim \mathcal{N}(0,\, 1) \\
  \alpha_{\text{duracion}} &\sim \mathcal{N}(0,\, 1) \\
  \alpha_{\text{estacion}} &\sim \mathcal{N}(0,\, 1) \\
  \alpha_{\text{protocolo}} &\sim \mathcal{N}(0,\, 1)
\end{align}

\medskip
\textbf{\textcolor{ghostgray}{Eje Latente de Acondicionamiento ($\omega$):}}
\begin{equation}
  \omega \sim \mathcal{PG}(b=1,\; c=0)
\end{equation}

\begin{notabox}
\small\textit{Nota teórica:} La distribución Pólya-Gamma (Polson et al., 2013) requiere dos propiedades: el parámetro de sucesos se fija en $b=1$ al estar condicionado a la verosimilitud binaria de eBird, y $c=0$ ancla la media en el centro, comportándose como un acondicionador neutral universal.
\end{notabox}

\subsection{Justificación Metodológica de la Prior No Informativa}

El uso estricto de una prior neutra $\mathcal{N}(0, 1)$ se fundamenta en una \textbf{premisa de objetividad por falta de experticia a priori}:

\begin{itemize}
  \item Como equipo, no somos ornitólogos expertos como para asegurar puntualmente qué tanto un microgramo de Ozono diezma a las poblaciones de \textit{Z. capensis}. Al no saberlo, nos declaramos ignorantes fijando la expectativa central en cero ($\mu_0 = 0$).
  \item Modificar la media o la varianza para ``ajustar'' el modelo, aunque fuera en un mísero $0.01$, implicaría insertar conocimiento preconcebido, sesgando irreparablemente el diseño.
  \item Si el modelo decreta que la polución impacta la ocupación, ese descubrimiento proviene \textbf{única y exclusivamente} de la presión de los datos silentes de eBird sobre la hoja en blanco empírica, no de un juicio personal embebido en el código.
\end{itemize}

\hrule
\vspace{1em}

% ============================================================
\section{La Likelihood (Verosimilitud)}
% ============================================================

\subsection{Separación de Parámetros: Beta vs. Alpha}

En los modelos jerárquicos de ocupación se independiza matemáticamente el ecosistema real del error humano de observación. El modelo se fractura en dos mundos que actúan como válvulas:

\begin{itemize}
  \item \textbf{\textcolor{ecogreen}{Mundo Ecológico ($\beta$):}} Modelan la biología de la especie (ocupación latente, $\psi$). Responden si el ecosistema permite que el ave viva en el sitio.
  \item \textbf{\textcolor{humanred}{Mundo Humano ($\alpha$):}} Modelan el error observacional (detección, $p$). Responden si el humano registró al ave, asumiendo que esta sí estaba presente.
\end{itemize}

\subsection{La Likelihood Aumentada y el Origen de la Pólya-Gamma}

El proceso que deriva la verosimilitud final sigue una lógica de tres escalones:

\medskip
\textbf{Paso 1: La Naturaleza Binaria (Bernoulli)}

En cada evento de observación, el resultado es binario: ``vi al Copetón (1)'' o ``no lo vi (0)''. La verosimilitud de un solo dato $y_i$ (presencia/ausencia) para un sitio $i$ se define como:

\begin{equation}
  P(y_i \mid \psi_i) = \psi_i^{\,y_i}\,(1 - \psi_i)^{1-y_i}
\end{equation}

\medskip
\textbf{Paso 2: El Enlace Logístico (Inverse-Logit)}

Para vincular la probabilidad biológica $\psi$ con las variables ambientales $X_i\beta$, se usa la función Inverse-Logit, que fuerza cualquier valor real al intervalo $(0,1)$:

\begin{equation}
  \psi_i = \frac{e^{X_i\beta}}{1 + e^{X_i\beta}}
\end{equation}

Al sustituir esta expresión en la Bernoulli del Paso~1, se llega a la expresión computacionalmente difícil:

\begin{equation}
  P(y_i \mid X_i\beta) = \frac{\left(e^{X_i\beta}\right)^{y_i}}{1 + e^{X_i\beta}}
\end{equation}

\begin{notabox}
\small \textbf{¿Por qué es difícil?} El denominador $\left(1 + e^{X_i\beta}\right)$ es no lineal. Al multiplicarlo con la Prior Normal, los exponentes no se cancelan limpiamente, por lo que no existe una solución analítica directa (``solución cerrada''). El algoritmo se vería obligado a adivinar la forma de la curva a ciegas, iteración a iteración.
\end{notabox}

\medskip
\textbf{Paso 3: El Acondicionamiento Estocástico (Identidad de Pólya-Gamma)}

Para romper ese denominador intratable, Clark (2019) aplica la identidad de Polson et al.\ (2013). Esta identidad permite reescribir la probabilidad como una integral ponderada por la variable latente $\omega$:

\begin{equation}
  \underbrace{\dfrac{\left(e^{\psi}\right)^y}{1+e^{\psi}}}_{\text{Inverse-Logit}}
  \;=\;
  \underbrace{\dfrac{1}{2}\exp\!\left((y - \tfrac{1}{2})\psi\right)}_{\text{Capa Lineal}}
  \int_0^{\infty}
  \underbrace{\exp\!\left(-\dfrac{\omega\,\psi^2}{2}\right)}_{\text{Capa Cuadrática (Normal)}}
  p(\omega \mid 1,\, 0)\; d\omega
  \label{eq:pg_identity}
\end{equation}

De la ecuación~\eqref{eq:pg_identity} se desprenden dos consecuencias clave:

\begin{itemize}
  \item \textbf{El Transformador Latente ($\omega$):} No es un parámetro fijo, sino una \textbf{variable latente estocástica}. En cada iteración del modelo, se le pregunta a la distribución $\mathcal{PG}(1,0)$: \textit{``¿Qué peso necesito para que este dato de eBird se vea lineal?''}. $\omega$ responde con un número que ``endereza'' localmente la curva del Logit.
  \item \textbf{La Linealización:} El término dentro de la integral tiene la forma $\exp\!\left(-\tfrac{\omega\psi^2}{2}\right)$, que es una \textbf{estructura cuadrática en el exponente}. En estadística, una parábola en el exponente es la firma matemática de una \textbf{Distribución Normal}.
\end{itemize}

\begin{notabox}
\small \textbf{¿Por qué esto es relevante?} Porque la Prior también es Normal. Al tener ahora una Likelihood que se comporta como Normal y una Prior Normal, el producto de ambas produce una \textbf{solución cerrada} (otra Normal). El modelo deja de ``adivinar'' y puede calcular la respuesta exacta de forma directa.
\end{notabox}

La \textbf{Likelihood Aumentada de Datos Completos} para los $N$ sitios y $K_i$ visitas por sitio queda entonces como:

\begin{equation}
  L\!\left(y, z, \omega \mid \beta, \alpha\right)
  = \prod_{i=1}^N
    \underbrace{P\!\left(z_i,\, \omega^\psi_i \mid \beta\right)}_{\text{Capa Biológica}}
    \cdot
    \prod_{j=1}^{K_i}
    \underbrace{P\!\left(y_{i,j},\, \omega^p_{i,j} \mid z_i,\, \alpha\right)}_{\text{Capa Humana}}
\end{equation}

Donde:
\begin{itemize}
  \item \textbf{Válvula Biológica:} $z_i \sim \text{Bernoulli}(\psi_i)$. La contaminación intenta fijar $\psi_i \to 0$.
  \item \textbf{Válvula de Observación:} $y_{i,j} \sim \text{Bernoulli}(z_i \cdot p_{i,j})$. El candado biológico maestro: si $z_i = 0$, todo el esfuerzo humano se multiplica por cero ($0 \times p = 0$). No se puede detectar lo que no está presente.
\end{itemize}

\hrule
\vspace{1em}

% ============================================================
\section{Construcción de la Posterior Paso a Paso}
% ============================================================

\subsection{La Fórmula de la Posterior Conjunta}

Aplicando la regla de Bayes al modelo aumentado, la \textbf{Posterior Conjunta} es proporcional al producto de la Likelihood Aumentada por todas las Priors:

\begin{equation}
  \boxed{
  P\!\left(\beta, \alpha, \mathbf{z}, \boldsymbol{\omega} \mid \mathbf{y}\right)
  \;\propto\;
  \underbrace{L\!\left(\mathbf{y}, \mathbf{z}, \boldsymbol{\omega} \mid \beta, \alpha\right)}_{\text{Datos (Likelihood)}}
  \times
  \underbrace{\pi(\beta)\;\pi(\alpha)\;\pi(\boldsymbol{\omega})}_{\text{Ignorancia (Priors)}}
  }
\end{equation}

Al sustituir los componentes se observa la fusión total:

\begin{multline}
  P\!\left(\beta, \alpha, \mathbf{z}, \boldsymbol{\omega} \mid \mathbf{y}\right)
  \;\propto\;
  \prod_{i=1}^N
  \left[
    P\!\left(z_i, \omega^\psi_i \mid \beta\right)
    \prod_{j=1}^{K_i}
    P\!\left(y_{i,j}, \omega^p_{i,j} \mid z_i, \alpha\right)
  \right] \\
  \times\;
  \mathcal{N}(\beta \mid 0, 1)
  \;\times\;
  \mathcal{N}(\alpha \mid 0, 1)
  \;\times\;
  \mathcal{PG}(\omega \mid 1, 0)
\end{multline}

\subsection{La Solución por Bloques: Gibbs Sampling}

Como no es posible resolver la ecuación conjunta de una sola vez, el modelo la rompe en \textbf{4 bloques iterativos} que se actualizan en ciclo (Muestreo de Gibbs):

\begin{enumerate}[label=\textbf{Bloque \arabic*.}]
  \item \textbf{Estado Latente ($z$):} Se imputa el estado real del ave en los sitios donde no hubo detección, comparando la probabilidad de extirpación ecológica contra la de fallo observacional.
  \item \textbf{Acondicionamiento ($\omega$):} Se generan las variables fantasma Pólya-Gamma que aplanan el Logit para el siguiente paso.
  \item \textbf{Ecología ($\beta$):} Se calculan los efectos de $\text{PM}_{10}$ y $\text{O}_3$. Gracias a $\omega$, este paso tiene ahora una distribución \textbf{Normal} de solución directa.
  \item \textbf{Humano ($\alpha$):} Se calculan los efectos del esfuerzo del pajarero. Igualmente, resulta en una distribución \textbf{Normal} de solución directa.
\end{enumerate}

\subsection{La Mecánica de la Solución Cerrada (Conjugación)}

En cada iteración del Bloque 3 y 4, la Posterior de los coeficientes se calcula mediante conjugación \textbf{Normal--Normal}. Las fórmulas de actualización son:

\medskip
\textbf{1. Varianza Posterior Actualizada ($V_n$):}
\begin{equation}
  V_n = \left(X^T \Omega\, X + V_0^{-1}\right)^{-1}
\end{equation}

\noindent donde $\Omega = \text{diag}(\omega_1, \dots, \omega_N)$ es la matriz diagonal de los pesos fantasmas y $V_0^{-1}$ es la precisión de la Prior (para $\mathcal{N}(0,1)$, $V_0^{-1} = I$).

\medskip
\textbf{2. Media Posterior Actualizada ($\mu_n$):}
\begin{equation}
  \mu_n = V_n \left(X^T \kappa + V_0^{-1}\,\mu_0\right)
\end{equation}

\noindent donde $\kappa_i = y_i - \tfrac{1}{2}$ son los datos de eBird centrados y $\mu_0 = \mathbf{0}$ es la media inicial de ignorancia.

\begin{notabox}
\small Este mecanismo convierte un problema de optimización no lineal y lento en una serie de \textbf{operaciones algebraicas directas y exactas}. Es por esto que los modelos de ocupación con augmentación Pólya-Gamma son órdenes de magnitud más rápidos que los algoritmos Metropolis-Hastings anteriores.
\end{notabox}

\subsection{La Convergencia de las Cadenas de Markov}

El ciclo de 4 bloques se repite miles de veces. Al principio, los valores saltan erráticamente (período de \textit{burn-in}, ${\sim}1{,}000$ iteraciones, se descartan). Tras la convergencia, los siguientes ${\sim}10{,}000$ valores constituyen la muestra empírica de la Posterior: la respuesta final del modelo.

\hrule
\vspace{1em}

% ============================================================
\section{Resultado Final: Los Parámetros Posteriores}
% ============================================================

Al terminar el proceso, el modelo entrega miles de muestras de la Posterior. Para cada contaminante se obtienen:

\begin{enumerate}
  \item \textbf{Estimación puntual:} La media o mediana de la distribución posterior.
  \item \textbf{Incertidumbre cuantificada:} Intervalos de Credibilidad al 95\%.
\end{enumerate}

\medskip
\textbf{Interpretación Biológica:}
Si el intervalo de credibilidad al 95\% del coeficiente $\beta_{\text{PM10}}$ no incluye al cero y es completamente negativo:

\begin{equation}
  \text{IC}_{95\%}[\beta_{\text{PM10}}] < 0
\end{equation}

\noindent habremos demostrado que el $\text{PM}_{10}$ reduce significativamente la probabilidad de ocupación del Copetón (\textit{Zonotrichia capensis}) en Bogotá, \textbf{habiendo controlado todo el ruido de imperfección de detección humana} a través de la capa jerárquica de los parámetros $\alpha$.

\vspace{1em}
\hrule
\vspace{0.5em}
\begin{center}
  \small\textit{Referencia principal: Clark \& Altwegg (2019). Efficient Bayesian analysis of occupancy models with logit link functions. \textbf{Ecology and Evolution}. \\ Técnica de aumento: Polson, Scott \& Windle (2013). \textbf{Journal of the American Statistical Association}.}
\end{center}

\end{document}
